\documentclass{article}

\textwidth=6.5in
\textheight=8.5in
\voffset=-54pt
\oddsidemargin=0in
\evensidemargin=0in
\setlength{\parskip}{8pt}
\setlength{\parindent}{0pt}

\usepackage{amsmath, amssymb}
\usepackage{ifthen}

\usepackage[inline]{enumitem}
\setlist{topsep=0pt, parsep=8pt}

\usepackage[rgb,table]{xcolor}\convertcolorsUfalse
\usepackage{graphicx}

% change hyperref options
\PassOptionsToPackage{hyphens}{url}
\usepackage[bookmarks,colorlinks,linkcolor=black,citecolor=blue,pdfstartview=FitH,urlcolor=blue]{hyperref}
\hypersetup{pdfpagemode=UseNone}
\usepackage{bookmark}

\usepackage{graphicx}
\usepackage{multicol}
\usepackage{framed}
\usepackage{array}

\usepackage{icomma}

\usepackage{marginnote}
\usepackage[colorinlistoftodos]{todonotes}
\renewcommand{\marginpar}{\marginnote}

\usepackage{soul}

\usepackage{pgfplots}
\pgfplotsset{compat=1.18}  
\pgfplotscreateplotcyclelist{mycolor}{black, blue, red, green!70!black, magenta, purple, cyan, orange }  
\pgfplotsset{
  disabledatascaling,
  cycle list=tikzcolor, 
  axis lines=middle,
  every axis plot/.style={no marks, thick},
  every axis/.style={
    enlarge x limits=false,
    enlarge y limits=auto,
    xlabel style={at={(current axis.right of origin)},anchor=west},
    ylabel style={at={(current axis.above origin)},anchor=south},
    % extra x and y ticks for asymptotes
    extra x tick style={grid=major, grid style={dashed,black}},
    extra y tick style={grid=major, grid style={dashed,black}},
    /pgf/number format/1000 sep={},
  },    
  tick label style = {font=\footnotesize, fill=\currentbgcolor, inner sep=0pt},    
  xticklabel shift=1pt,    
  scaled ticks=false,
  scaled y ticks=false,
  unbounded coords=jump,
  samples=101,
  minor tick length=0,
  scatter/classes={
    open={mark=*,draw=black,fill=white, mark size=2, thin},
    closed={mark=*,draw=black,fill=black, mark size=2, thin}
  },
  width=3in,
}
\pgfplotsset{open/.style={only marks, mark=*, fill=white, thin, mark size=2}}
\pgfplotsset{closed/.style={only marks, mark=*, draw=black, fill=black,
    thin, mark size=2}}
\pgfplotsset{labeled/.style={nodes near coords, only marks, mark=*, draw=black, fill=black, thin, mark size=2, point meta=explicit symbolic}}

\pgfplotsset{gridgraph/.style={clip=false, axis line style={shorten
      >=-8pt}, xlabel style={xshift=8pt}, ylabel style={yshift=8pt},
    enlargelimits=false, grid=major, tick style={black}}} 

\usepackage{tikz}
\usetikzlibrary{
  matrix,%
  % enables the snake paths
  decorations.pathmorphing,%
  % enables bigger arrows
  decorations.markings,%
  decorations.pathreplacing,%
  decorations.text,%
  calc,%
  patterns,%
  shapes.geometric,%
  arrows,
  decorations.shapes,
  positioning,
  plotmarks,
  intersections
}
\tikzset{>=latex} % sets default arrow type in tikz
\tikzset{font=\normalsize}
\tikzset{mark size=1.5pt, mark options=thin}
\tikzset{pin distance=4pt,
  every pin edge/.style={<-, thin, shorten <= -2pt}}

\interfootnotelinepenalty=10000
\renewcommand{\thefootnote}{(\arabic{footnote})}

\usepackage{multirow, bigdelim}

% change to \usepackage[sols]{handout} to see solutions
\usepackage[sols, grading, workshop]{handout}

\newcommand\iscurrentenumi[1]{%
  \ifthenelse{\equal{#1}{\theenumi}}{}{#1}}
\setenumerate[1]{ref=\#\arabic*}
\setenumerate[2]{ref=\protect\iscurrentenumi{\theenumi}(\alph*)}
\newcommand\enumiiprefix[2]{%
  \ifthenelse{{\equal{#1}{\theenumi}}\AND{\equal{#2}{(\alph{enumii})}}}{}{}%
  \ifthenelse{{\equal{#1}{\theenumi}}\AND{\NOT{\equal{#2}{(\alph{enumii})}}}}{#2}{}%
  \ifthenelse{\equal{#1}{\theenumi}}{}{#1#2}%
}
\setenumerate[3]{ref=\protect\enumiiprefix{\theenumi}{(\alph{enumii})}\roman*}

\makeatletter

% underline links               
\let\@href=\href
\renewcommand{\href}[2]{\@href{#1}{\ul{#2}}}

\newcommand\calculator{
  \begin{tikzpicture}[x=0.15ex, y=0.15ex, baseline=1.5pt]
    \begin{scope}[thick]
      \draw (0, 0) -- (10, 0) -- (10, 14) -- (0, 14) -- cycle;
      \foreach \x in {10/3, 20/3}
      {
        \draw (\x, 0) -- (\x, 10);
      }
      \foreach \y in {10/3, 20/3, 10}
      {
        \draw (0, \y) -- (10, \y);
      }
    \end{scope}
  \end{tikzpicture}
}
\newcommand{\calcitem}{\refstepcounter{enum\romannumeral\@enumdepth}%
\item[\calculator \csname labelenum\romannumeral\@enumdepth\endcsname]}

\makeatother

\definecolor{purple}{rgb}{0.5,0,1.0}
\definecolor{hotpink}{rgb}{1, 0, 0.5}
\definecolor{skyblue}{rgb}{0, 0.5, 1}

\newcommand{\goal}[1]{\textbf{\textcolor{skyblue}{#1}}}
\newcommand{\indvar}[1]{\textbf{\textcolor{hotpink}{#1}}}


\newcommand{\psreflection}[2][]{

  \ifthenelse{\not\equal{#1}{nocite}}{
    \begin{enumerate}[resume]
    \item
      \begin{problem}[1in]
        Please cite any help you got on this problem set:
        \begin{itemize}[nosep]
        \item If you worked with other students, please list their names here.
        \item If you used resources other than those on our Canvas site, please list them here.
        \end{itemize}
      \end{problem}

    \end{enumerate}
  }

  \probonly{%
    #2

    \newpage

    \emph{Reflection space.} It's a good habit to take a few minutes at the end of each pset to reflect on what you learned and jot down notes to your future self. Here are a few reflection questions to get you started. Nothing you write here will be graded; it's just for you!
    \begin{itemize}
    \item Take a few minutes to look back at the problems. How could you explain to someone else how to do each problem? What was your strategy, and how did you know to use that strategy? If there are problems where you don't feel able to answer this, write down which ones they are. Come back to them in a few days when we post the homework solutions!

      \vspace{1.5in}

    \item Take a look back at the learning objectives on today's worksheet; how are you feeling about each one? If there are ones you know you need more practice with, write those down here.

      \vspace{1.5in}

    \item What did you find most challenging about this material? Do you have any lingering questions about it? If so, write them down here so you don't forget them. At some point, stop by office hours or MQC to ask!

      \vfill
    \end{itemize}      
  }
}

\usepackage{fancyhdr}
\lhead{\textcolor{white}{This content is protected and may not be shared, uploaded, or distributed.}}
\rhead{}
\rfoot{\vspace*{\fill}\footnotesize\textcopyright{} \emph{Harvard Math Ma}}
\renewcommand{\headrulewidth}{0pt}
\pagestyle{fancy}
\begin{document}

\begin{center}
  \large\textsc{Problem Set 7 - Problem-solving day}
\end{center}

\probonly{There is no Edfinity for this problem set!}

\begin{enumerate}
\item  Tom and Harry are both trying to graph $f(x) = \frac{5}{x} -2$. Their approach is to use graph transformations.
  \begin{itemize} 
  \item Harry stretches the graph of $g(x)=\frac{1}{x}$ vertically by a factor of 5 and then shifts that down 2 units.
  \item Tom shifts the graph of $g(x) = \frac{1}{x}$ down 2 units and then stretches that vertically by a factor of 5.
  \end{itemize}   

  \begin{enumerate} 
  \item

    \begin{grading}
      3 points: 1 for each graph and 1 for how they're different.
    \end{grading}

    \begin{problem}[\vfill]
      Draw the graphs that Harry and Tom end up with. Explain how they're different.
    \end{problem}

    \begin{solution} 
      This is the graph that Harry ended up with. Note that Harry's graph has a horizontal asymptote at $y=-2$.
      \begin{center}
        \begin{tikzpicture}
          \begin{axis}[ymin=-12, ymax=12, ytick={-2}, xtick=\empty]
            \addplot+[domain=-10:-.35] {(5/x) -2};
            \addplot+[domain=.35:10] {(5/x) -2};
            \addplot+[dashed, domain=-10:10] {-2}; 
          \end{axis}
        \end{tikzpicture}
      \end{center}
      This is the graph that Tom ended up with. His graph has a horizontal asymptote at $y=-10$.
      \begin{center}
        \begin{tikzpicture}
          \begin{axis}[ymin=-20, ymax=4, ytick={-10}, xtick=\empty]
            \addplot+[domain=-10:-.35] {5*(1/x -2)};
            \addplot+[domain=.35:10] {5*(1/x -2)};
            \addplot+[dashed, domain=-10:10] {-10}; 
          \end{axis}
        \end{tikzpicture}
      \end{center}
    \end{solution} 

  \item

    \begin{grading}
      1 point
    \end{grading}

    \begin{problem}[\stretch{0.25}]
      Who ends up with the correct graph and why?
    \end{problem}

    \begin{solution} 
      Harry's graph is correct. We can tell because the graph of $f$ should have a horizontal asymptote at $y=-2$.
      The order we should perform a vertical shift and a vertical stretch follows the order of operations.
    \end{solution}

  \item

    \begin{grading}
      1 point
    \end{grading}

    \begin{problem}[\nstretch{0.25}]
      What is the algebraic formula for the incorrect graph?
    \end{problem}

    \begin{solution} 
      $f(x) = 5(\frac{1}{x} - 2) $.  
    \end{solution}

  \end{enumerate} 

\item  

  Mono Lake is a a large lake in California. Because it has no outlets, high levels of salts accumulate in the lake.

  \href{https://www.monolake.org/learn/stateofthelake/}{According to
    the Mono Lake Committee}, ``In 1941, the Los Angeles Department of
  Water \& Power (DWP) began diverting water from Mono Lake’s
  tributary streams, sending it 350 miles south to meet the growing
  water demands of Los Angeles.

  As a result, over the next 40 years Mono Lake dropped by 45 vertical
  feet, lost half its volume, and doubled in salinity—threatening the
  survival of the nesting California Gull population, air quality with
  toxic dust storms, and this unique and critical ecosystem.''

  Let $f(w)$ be the salinity (in grams of salt per liter of water) when the water elevation is $w$ feet above sea level, and let $g(t)$ be the water elevation (in feet above seat level) $t$ years after 1900. Here's a graph of $f$, based on \href{https://www.waterboards.ca.gov/waterrights/water_issues/programs/mono_lake/docs/backgroundinfo.pdf}{data from the California Water Boards}:

  \begin{tikzpicture}
    \begin{axis}[xlabel={water elevation (feet above sea level)}, xlabel style={at=(ticklabel cs:0.5), anchor=north, yshift=-4pt}, ylabel={salinity (g/L)}, ylabel style={at=(ticklabel cs:0.5), anchor=south, align=center, rotate=90}, ymin=0, xmax=6420, disabledatascaling=false, gridgraph, width=3.5in, height=2.25in, ymax=120, ytick={20, 40, ..., 120}]
      \addplot+[domain=6370:6420] {4673.282384550384/(-6324.792447354172 + x)};
    \end{axis}
  \end{tikzpicture}

  And here's a graph of $g$, using
  \href{https://www.monobasinresearch.org/data/levelyearly.php}{data
    from the Mono Lake Committee}:

  \begin{tikzpicture}
    \begin{axis}[ylabel=water elevation (feet above sea level), ylabel style={at=(ticklabel cs:0.5), anchor=south, xshift=-4pt, text width=2.5in, align=center, rotate=90}, ticks=major, grid=major, axis y discontinuity=parallel, width=4.5in, height=2.75in, disabledatascaling=false, xlabel={years after 1900}, xlabel style={at=(ticklabel cs:0.5), anchor=north, yshift=-4pt}, ymin=6360, ytick={6370, 6380, ..., 6430}]
      \addplot[smooth] file {data/mono-lake-water-level.txt};
    \end{axis}
  \end{tikzpicture}

  \pspace{\newpage}

  \begin{enumerate}
  \item

    \begin{grading}
      1 point. The key thing they should identify is that the graph of $f$ shows salinity as a function of water elevation, not time. 
    \end{grading}

    \begin{problem}[\vfill]
      Your friend says, ``I'm confused! The salinity of Mono Lake doubled between 1941 and 1981, but the graph of $f$ (salinity) is decreasing. How can that be?'' Write an explanation for your friend. 
    \end{problem}

    \begin{solution}
      Here's a possible explanation.

      On the graph, the independent variable is water elevation, not time. It shows that salinity decreases as water
      elevation increases, not as time passes.
    \end{solution}

  \item

    \begin{grading}
      1 point
    \end{grading}

    \begin{problem}[\stretch{0.5}]
      Write an expression in function notation that gives the salinity of Mono Lake in 1941.
    \end{problem}

    \begin{solution}
      $f(g(41))$
    \end{solution}

  \calcitem

    \begin{grading}
      2 points: 1 for the estimate (which may be slightly different from the estimate in the solutions, since they need to eyeball things), 1 for units
    \end{grading}

    \begin{problem}[\vfill]
      Use the graphs to estimate the average rate of change of salinity with respect to time between 1940 and 1980; include units in your answer.
    \end{problem}

    \begin{solution}
        Using the definition of average rate of change,
        \[
        \parbox{4cm}{\centering average rate of change\\of \textbf{salinity} \\with respect to \textbf{time}} = 
        \hspace{1em}
        \frac{\text{change in salinity}}{\text{change in time}}.
        \]
        The change in time is $40$ years. To estimate the change in salinity, we can use the graphs of $f$ and $g$
        to estimate $f(g(40))$ and $f(g(80))$. First, $g(40) \approx 6417$ and $g(80) \approx 6373$. Using these values
        $f(6417) \approx 52$ and $f(6374) \approx 94$. Thus, the change in salinity, $f(g(80))-f(g(40))$,
        is approximately 42 grams per liter.

        Plugging in these values, 
        \[
        \parbox{4cm}{\centering average rate of change\\of \textbf{salinity} \\with respect to \textbf{time}} \approx 
        \hspace{1em}
        \frac{42 \text{ grams per liter}}{40 \text{ years}} = 1.05 \text{ grams per liter per year}
        \]
    \end{solution}

  \calcitem

    \begin{grading}
      1 point: 0.5 for the estimate, 0.5 for units
    \end{grading}

    \begin{problem}[\vfill\newpage]
      Use the graphs to estimate the average rate of change of water elevation with respect to time between 1940 and 1980; include units in your answer.
    \end{problem}

    \begin{solution}
       Using the definition of average rate of change,
        \[
        \parbox{4cm}{\centering average rate of change\\of \textbf{water elevation} \\with respect to \textbf{time}} = 
        \hspace{1em}
        \frac{\text{change in water elevation}}{\text{change in time}}.
        \]
        Again, the change in time is $40$ years. We estimated before that $g(3) \approx 6417$ and $g(80) \approx 6373$,
        so the change in water elevation, $g(80)-g(40)$, is approximately $-44$ feet.

        Plugging in these values,
        \[
        \parbox{4cm}{\centering average rate of change\\of \textbf{water elevation} \\with respect to \textbf{time}} 
        \approx \hspace{1em}
        \frac{-44 \text{ feet}}{40 \text{ years}} = -1.1 \text{ feet per year}
        \]
        
    \end{solution}
  \end{enumerate}

\item \label{fall2015-1a-gateway-overview-25} \label{modeling-with-functions-can-1}

  \begin{grading}
    8.5 points:
    \begin{itemize}[nosep]
    \item 1 for drawing a picture
    \item 0.5 for labeling the picture appropriately
    \item 1 for stating the goal
    \item 2 for expressing cost in terms of radius and height: 1 for correctly relating the cost to areas, 1 for correctly finding the relevant areas
    \item 1 for writing a correct equation about volume
    \item 1 for solving their volume equation for height in terms of the radius 
    \item 1 for plugging that into the cost
    \item 1 for the domain 
    \end{itemize}
  \end{grading}

  \begin{problem}[\newpage]
    An engineer is designing an aluminum soft drink can, which will have a volume of $128 \pi$ cubic centimeters. The can is to be made out of two types of aluminum. The top and bottom of the can need to be extra durable, so they will be made out of aluminum costing 7 cents per square cm. The rest of the can will be made out of thinner aluminum costing 4 cents per square cm. Express the cost of the can in terms of its radius. What is the domain of your function?

    Whenever you do a modeling problem like this, you should follow the guidelines in \href{https://people.math.harvard.edu/~jjchen/mathma/docs/homework/PS06.html}{Problem Set 6}:
    \begin{itemize}[nosep]
    \item State the goal and the information you're given.
    \item Draw a picture, explicitly labeling both variables and constants.
    \item Clearly communicate your thought process, one step at a time. 
    \end{itemize}
  \end{problem}

  \begin{solution}
    \hl{\textbf{The goal}} is to express the \goal{total cost} of the can in terms of its \indvar{radius}.
    \hl{\textbf{We know}} that the volume of the can is $128\pi$ cubic centimeters, 
    the cost per area of aluminum used for the top and bottom of the can is 7 cents per square centimeter, 
    and the cost per area of the aluminum used for the sides of the can is 4 cents per square centimeter. 

    We can \hl{\textbf{draw a picture}} to get a better sense of what's going on.
    \begin{center}
      \begin{tikzpicture}
        \draw (0,0) ellipse (1.25 and 0.5);
        \draw[dashed] (0, 0) -- node[midway, above] {$r$} (1.25, 0); 
        \draw (-1.25,0) -- (-1.25,-3.5);
        \draw (-1.25,-3.5) arc (180:360:1.25 and 0.5);
        \draw [dashed] (-1.25,-3.5) arc (180:360:1.25 and -0.5);
        \draw (1.25,-3.5) -- (1.25,0);
        \draw[|-|] (-1.75, -3.5) -- node[midway, inner sep=6pt, fill=white] {$h$} (-1.75,0); 
      \end{tikzpicture}
    \end{center}
    Now, let's \hl{\textbf{write a formula}} for the cost of the can.
    \begin{align*} 
      \text{ cost of can } &= \text{cost of top and bottom } + \text{
        cost of sides } \\
      &= (\text{7 cents/square cm}) (\text{area of top and
        bottom}) + (\text{4 cents/square cm}) (\text{area of sides}) \\
      \intertext{The top and bottom are disks of radius $r$, so they
        each have area $\pi r^2$.  The area of the sides is $2 \pi r h$, 
        where $h$ is the height of the can.}
      &= 7 (\pi r^2 + \pi r^2) + 4 (2 \pi r h) \\
      &= 14 \pi r^2 + 8 \pi r h. \\
      \intertext{To write the cost in terms of $r$, we need to \hl{\textbf{relate the variables}} by
      \hl{\textbf{using the information we know.}} We want to 
        write $h$ in terms of $r$.  We were told that the volume of
        the can is to be $128 \pi$, and the volume of a cylinder is $V= \pi r^2 h$. Solving $128 \pi = \pi r^2 h$ for
        $h$ in terms of $r$ gives $h= \frac{128}{r^2}$. Substituting this into our formula for the cost, we get}
      &= 14 \pi r^2 + 8 \pi r \left( \frac{128}{r^2} \right) \\
      &= \boxed{14 \pi r^2 + \frac{1024 \pi}{r}}.
    \end{align*}
    The domain of this function represents all possible values of $r$. All we can say is that $r$ must be positive (our can could be extremely skinny and very tall, or very wide and very short; it's still possible to have a volume of $128 \pi$ cubic centimeters). So the domain is $\boxed{r > 0}$ or $\boxed{(0, \infty)}$. 
  \end{solution}

\end{enumerate}
\psreflection{For next class, read \S 4.1 and Example 4.8 of \S 4.4.}
\end{document}
